\subsection{Overall Framework}
\label{sec:framework}

The proposed CLIP-Enhanced Semantic Manifold Diffusion (CESM-Diff) framework is designed to address the challenges of zero-shot fault diagnosis by bridging the high-dimensional feature manifold of industrial sensor signals with the pre-trained, context-rich semantic manifold of the Contrastive Language-Image Pretraining (CLIP) text encoder. As shown in Figure 3, the overall architecture consists of four integrated components: (1) a multi-modal feature extractor that processes diverse temporal sensor inputs; (2) a CLIP-driven semantic embedding layer that generates high-fidelity textual descriptors; (3) a Dual-Path Semantic Diffusion module that synchronizes the generative processes in both feature and semantic spaces; and (4) a Semantic Manifold Interpolation module that "hallucinates" the features for unseen fault categories.

\subsubsection{Workflow Overview}
The system operates in a two-stage training paradigm to ensure both reconstruction fidelity and semantic grounding. During the seen-class training phase, the model learns a robust mapping between the vibration patterns $\mathbf{x} \in \mathcal{X}_{seen}$ and their corresponding textual labels $\mathcal{L}_{seen}$. This mapping is established within a collective latent space $\mathcal{Z}$, where individual fault classes form distinct clusters according to their semantic descriptions, drawing upon principles from self-supervised representation learning~\cite{bardes2021vicreg,caron2021dino,chen2021mocov3}. The CLIP text encoder provides the "semantic anchor" for each cluster, while the diffusion-based generator $G$ learns to model the complex, non-linear distribution of the vibration features.

Once the manifold structure is established for the seen classes, the system enables zero-shot inference for novel fault categories $\mathcal{L}_{unseen}$. Specifically, the model retrieves the semantic embeddings $\mathbf{s}_{unseen}$ for the unseen classes from CLIP and utilizes the Semantic Manifold Interpolation (SMI) to navigate the latent trajectory between related seen faults. This allows for the "synthetic hallucination" of vibration features for classes never encountered during training. By sampling the reverse diffusion path conditioned on these interpolated embeddings, the generator $G$ produces a set of synthetic samples $\{\hat{\mathbf{x}}_{unseen}\}$ that correctly reside in the physical feature space. These synthetic samples are then used to train a simple, but effective, softmax classifier for the final diagnostic decision.

\subsubsection{Multi-Sensor Feature Extractor}
To handle the heterogeneity of industrial data, such as the multi-frequency signals in hydraulic systems or the multi-variable process data in the TEP dataset, we utilize a flexible convolutional feature extractor. For 1D vibration signals, the extractor is composed of multiple dilated convolutional layers that capture multi-scale temporal dependencies without a significant increase in parameter count. For 2D process data, we utilize a ResNet-like architecture~\cite{he2016deep} that treats the sensor-time matrix as an "image," capturing the spatial-temporal couplings between different process variables. The resulting feature vector $\mathbf{z}$ is then projected into the joint latent space $\mathcal{Z} \in \mathbb{R}^{512}$ to match the dimensionality of the CLIP embeddings. This ensures that the physical dimensions of the sensor data are mapped into a geometrically consistent structure where semantic relationships can be effectively exploited.

\subsubsection{CLIP-Driven Semantic Alignment}
The core innovation of our framework is the use of the CLIP text encoder as the primary source of semantic supervision. Unlike traditional zero-shot methods that rely on binary attributes or Word2Vec embeddings—which lack domain-specific technical context—CLIP has been pre-trained on a vast corpus of image-text pairs, allowing it to "understand" complex failure descriptions like "outer race pitting" or "reactor pressure oscillation." Each fault class $\mathcal{C}_i$ is represented by a descriptive prompt, such as "a vibration signal for a bearing with an inner race fault of 0.014 inch damage." The CLIP text encoder transforms these prompts into high-dimensional semantic anchors $\mathbf{s}_i$. The dual-path diffusion process then aligns the feature manifold $\mathcal{M}_{feat}$ with these anchors, ensuring that the geometric distance between two fault clusters in $\mathcal{Z}$ is proportional to their semantic similarity in the CLIP space. This alignment is reinforced by the Attribute Consistency Loss (ACL) discussed in Section 3.7 and the cross-path attention mechanism introduced in Section 3.5.

\subsubsection{Synergy and Information Exchange}
The interactive nature of the CESM-Diff framework is further enhanced by the bidirectional feed-back between the feature and semantic paths. At each diffusion step $t$, the Cross-Path Attention~\cite{vaswani2017attention} allows the feature extractor to "attend" to specific technical terms in the fault description, while the semantic refiner gains physical "grounding" from the noisy sensor data. This synergy ensures that the final diagnostic classifier is not only accurate but also interpretable, as its decisions are directly rooted in the semantic manifold derived from expert-curated fault descriptions. The overall system provides a scalable solution for zero-shot diagnosis, where new faults can be integrated simply by providing a descriptive textual prompt, eliminating the need for expensive data collection for every potential failure mode in industrial systems.
